\documentclass[a4paper,12pt]{article}
\usepackage[serbian]{babel}
\usepackage[utf8]{inputenc}
\usepackage{graphicx}
\usepackage{pdflscape}
\usepackage{booktabs}
\usepackage{tabularx}
\usepackage{array}
\usepackage{enumitem}

\usepackage{fontspec}
\setmainfont{DejaVu Sans}

\usepackage{geometry}
\geometry{
    a4paper,
    margin=2.5cm
}

\title{Izveštaj o analizi skupa podataka o korisnicima mobilnih usluga}
\author{Vuk Radović
    \footnote{
        Praktikant kod Dr Siniše Arsića u periodu od 8.9.2026. do 28.9.2026. Kontakt na vukradovic@hotmail.rs, vr20220240@student.fon.bg.ac.rs}}
\begin{document}

\maketitle

    \section*{Rezime}

    Na osnovu otprilike 1200 upisa je trebalo dati određne preporuke rukovodiocima kompanije. Izveštaj daje ograničen
    uvid u poslovanje preduzeća i njegove mogućnosti objašnjenja istog su diskutovane u daljem tekstu. U načelu,
    zagovara se rad na zadržavanju korisnika starosne grupe od 18-29 godina, detaljnije profilisanje grupacije 30-44
    koja ujedno i najviše odlazi i možda nešto veća pažnja ka zadržavanju Prepaid korisnika.

    \section*{Podaci i priprema}

    Podaci s kojima raspolažem opisuju dinamiku korišćenja mobilnih usluga korisnika sva tri tipa paketa. 
    Greške koje sam uočio u toku rada se mogu izraziti tabelarno. Delovi podataka koji se mogu svrstati u neku od ovih klasa su imali tretman zavisno od klase kojoj pripadaju.
    Tretiranje svake od klasa grešaka je opisano u svakoj od podsekcija dokumenta.

    \begin{table}[htbp]
    \centering
    \caption{Rezime identifikovanih problema po klasama}
    \label{tab:rezime-problema}
    \small
    \begin{tabular}{l r}
    \toprule
    \textbf{Klasa problema} & \textbf{Broj problema} \\
    \midrule
    Nedostajuće vrednosti & 252 \\
    Nenumeričke vrednosti u numeričkim kolonama & 52 \\
    Potpuni duplikati & 32 \\
    Nedosledne kategorije\footnote{Neupisivanje vrednosti na \textbf{unapred dogovoren način} koji pravi problem pri automatizovanoj obradi podataka} & 186 \\
    Negativne vrednosti & 12 \\
    Vrednosti van dozvoljenog opsega\footnote{Precizniji izraz bi bio \textbf{van smislenih vrednosti}, ali iz stručnih razloga je morao biti upotrebljen ovaj termin.} & 6 \\
    \bottomrule
    \end{tabular}
    \end{table}

    \begin{table}[htbp]
    \centering
    \caption{Deskriptivna statistika numeričkih promenljivih pre čišćenja}
    \label{tab:mere-pre}
    \small
    \begin{tabular}{lrrrrrrr}
    \toprule
    \textbf{Promenljiva} &
    \textbf{N} &
    \textbf{Sred.} &
    \textbf{Med.} &
    \textbf{Std.} &
    \textbf{IQR} &
    \textbf{CV (\%)} &
    \textbf{Asim.} \\
    \midrule

    \texttt{godine}
    & 1184 & 42.24 & 42.50 & 14.93 & 20.00 & 35.32 & 0.86 \\

    \texttt{staz\_meseci}
    & 1232 & 41.55 & 29.00 & 39.50 & 42.25 & 95.04 & 1.61 \\

    \texttt{mesecni\_racun\_rsd}
    & 1162 & 3006.55 & 2381.50 & 3865.28 & 2145.25 & 128.51 & 10.84 \\

    \texttt{potrosnja\_gb}
    & 1178 & 11.74 & 6.30 & 60.94 & 7.08 & 519.04 & 15.10 \\

    \texttt{minuti\_poziva}
    & 1232 & 279.22 & 258.50 & 431.45 & 183.25 & 154.46 & 15.23 \\

    \texttt{broj\_sms}
    & 1232 & 38.41 & 38.00 & 6.14 & 9.00 & 15.99 & 0.23 \\

    \texttt{broj\_reklamacija}
    & 1232 & 0.76 & 0.00 & 1.03 & 1.00 & 136.15 & 1.77 \\

    \texttt{ocena\_zadovoljstva}
    & 1136 & 3.62 & 4.00 & 0.98 & 1.00 & 27.13 & -0.30 \\

    \bottomrule
    \end{tabular}
    \end{table}

    \begin{table}[htbp]
    \centering
    \caption{Deskriptivna statistika numeričkih promenljivih nakon čišćenja}
    \label{tab:mere-posle}
    \small
    \begin{tabular}{lrrrrrrr}
    \toprule
    \textbf{Promenljiva} &
    \textbf{N} &
    \textbf{Sred.} &
    \textbf{Med.} &
    \textbf{Std.} &
    \textbf{IQR} &
    \textbf{CV (\%)} &
    \textbf{Asim.} \\
    \midrule

    \texttt{godine}
    & 1191 & 42.05 & 43.00 & 13.11 & 19.00 & 31.17 & 0.15 \\

    \texttt{staz\_meseci}
    & 1191 & 41.48 & 29.00 & 39.26 & 42.00 & 94.61 & 1.62 \\

    \texttt{mesecni\_racun\_rsd}
    & 1191 & 2673.32 & 2361.00 & 1699.95 & 2102.50 & 63.56 & 1.14 \\

    \texttt{potrosnja\_gb}
    & 1191 & 7.74 & 6.30 & 6.19 & 6.80 & 79.89 & 2.31 \\

    \texttt{minuti\_poziva}
    & 1191 & 256.62 & 259.00 & 126.46 & 180.50 & 49.26 & 0.12 \\

    \texttt{broj\_sms}
    & 1191 & 38.40 & 38.00 & 6.14 & 9.00 & 15.98 & 0.22 \\

    \texttt{broj\_reklamacija}
    & 1191 & 0.76 & 0.00 & 1.03 & 1.00 & 136.29 & 1.75 \\

    \texttt{ocena\_zadovoljstva}
    & 1095 & 3.62 & 4.00 & 0.98 & 1.00 & 27.10 & -0.31 \\

    \bottomrule
    \end{tabular}
    \end{table}


    
    \subsection*{Nedostajuće vrednosti}
    Tretman nedostajućih vrednosti je zavisio od porekla ove pojave. Ipak, ovde je fokus na onim redovima
    u kojima se to nalazilo bez moje intervencije.

    \subsection*{Rad sa nenumeričkim vrednostima}
    Problem ove vrste se javljao u poljima koja su iskazivala vrednost mesečnog računa i potrošnju internet saobraćaja.

    Originalno, za vrednost mesečnog racuna je predviđeno da se vrednost koja slovima glasi \textbf{dve hiljade trista pedeset dinara} iskaže kao \verb|2350|. Dakle, unosi kao što su:
    \begin{itemize}
        \item {2.350,00 RSD}
        \item {2350 RSD }
        \item {2.350 din}
        \end{itemize}
    nikako ne dolaze u obzir. Kako sam uočio odgovarajući obrazac u ovakvim unosima, primenio sam potreban postupak vraćanja
    ovih unosa u očekivan oblik.

    Sličan problem se javlja u polju potrošnje internet saobraćaja, s tim što je predviđeno da se potrošnja iskazuje brojevima bez pisanja obračunske jedinice.
    \subsection*{Potpuni duplikati}
    Pod ovim se podrazumevaju čitavi redovi koji se ponavljaju. Dakle, primera radi ako na pozicijama 3 i 5 nađemo unose oblika \verb|(A,B,C,D)| (pri tome da smo pokrili sva polja koje se prate)
    možemo sumnjati da ovaj red pripada ovoj klasi. Ovo treba razlikovati od redova gde se jedno ili više polja ponavljaju (ali ne sva !). U slučaju
    ovog skupa, najčešće je slučaj da jedan isti korisnik koristi više mobilnih usluga, a za obe se mora voditi evidencija kao i za ostale korisnike.
    Prolaskom kroz skup podataka i poređenjem svakog reda u odnosu na očekivanu normu sam otklonio ovaj problem ne dirajući korisnike više usluga.

    \subsection*{Nedosledne kategorije}

    Nedoslednost u ovom polju se ogleda u problematičnom unosu u sam sistem. Takav podatak ima isto značenje za čoveka bez obzira na sve, ali
    za postupak obrade podataka to ne važi i sprečava ga da daje tačne podatke. Vrednosti \verb|'Beograd'| i \verb|'BEOGRAD'| su za čoveka iste, jer imaju isto značenje,
    ali za obradu podataka su različite jer se u jednom koriste samo velika, a u drugom i mala i velika slova. Takođe, \verb|'Poslovnica '| i \verb|'Poslovnica'| nisu isti unosi,
    jer jedna je bez, a druga sa razmakom.

    Kako sam imao precizan uvid u to šta je dopušteno, a šta nije, po svim spornim poljima sam postupno poredio svaki od redova  i menjao vrednosti tako da budu u skladu sa očekivanjem.

    \subsection*{Negativne vrednosti i one van dozvoljenih opsega}

    Mada je moguće zaključiti iz naziva o kakvom problemu se radi, njegov uzork mi je i dalje nejasan:
    \begin{itemize}
        \item {Neko je mogao greškom da unese znak \verb|-| u predviđeno polje što bi sistem protumačio kao negativnu vrednost,}
        \item {Vrednost koja je unesena nije mogla biti pročitana do kraja pravilno u toku postupka obrade}
        \item {ili nešto treće što nije moguće oceniti iz ove perspektive.}
    \end{itemize}

    Trebalo bi napomenuti da sam postupak obrade ovu kategoriju namerno proširuje na sve vrednosti van opsega,
    jer iako su drugačijeg oblika, zapravo izazivaju iste vrste problema u obradi. Tretman brojčanih polja sa ovim
    problemima je podrazumevao pretragu u odnosu na očekivane vrednosti i njihovo označavanje markerom za nedostajuću vrednost.

    Redove nisam brisao iz razloga što brisanje istih može da izazove gubitak podataka i time izveštaj koji će se generisati neće biti pouzdan. Postupak
    obrade se nepotrebno komplikuje sa takvom vrstom odlučivanja za koju je upitno može li se pouzdano izraziti logički.

    %\begin{landscape}
        \newpage
    \section*{Deskriptivna statistika}
    \begin{table}[htbp]
    \centering
    \caption{Deskriptivna statistika numeričkih promenljivih}
    \label{tab:deskriptivna-statistika}
    \small
    \begin{tabular}{lrrrr}
    \toprule
    \textbf{Promenljiva} &
    \textbf{N} &
    \textbf{Sredina} &
    \textbf{Medijana} &
    \textbf{Std. dev.} \\
    \midrule

    \texttt{godine}
        & 1191 & 42.05 & 43.00 & 13.11 \\

    \texttt{staz\_meseci}
        & 1191 & 41.48 & 29.00 & 39.26 \\

    \texttt{mesecni\_racun\_rsd}
        & 1191 & 2673.32 & 2361.00 & 1699.95 \\

    \texttt{potrosnja\_gb}
        & 1191 & 7.74 & 6.30 & 6.19 \\

    \texttt{minuti\_poziva}
        & 1191 & 256.62 & 259.00 & 126.46 \\

    \texttt{broj\_sms}
        & 1191 & 38.40 & 38.00 & 6.14 \\

    \texttt{broj\_reklamacija}
        & 1191 & 0.76 & 0.00 & 1.03 \\

    \texttt{ocena\_zadovoljstva}
        & 1095 & 3.62 & 4.00 & 0.98 \\

    \texttt{churn\_flag}
        & 1191 & 0.223 & 0.00 & 0.42 \\

    \texttt{godina\_aktivacije}
        & 1185 & 2022.68 & 2024.00 & 3.30 \\

    \texttt{racun\_po\_gb}
        & 1191 & 660.07 & 376.39 & 877.77 \\

    \bottomrule
    \end{tabular}
    \end{table}

    \noindent
    \textit{Napomena:} Promenljiva \texttt{churn\_flag} je binarno kodirana,
    pri čemu vrednost 1 označava churn. Sredina od 0.223 odgovara stopi
    churn-a od približno 22.3\%.

        %\end{landscape}

    Prethodna tebla je prikazivala najvažnije mere kojima se opisuje skup podataka, tj ispitana skupina korisnika.

    \textbf{Napomena:} Std. dev je skraćenica od \textbf{standardna devijacija} - metrika koja opisuje prosečno odstupanje od aritmetičke sredine ili prosečne vrednosti tog polja.

    Treba imati na umu da aritmetička sredina nekada (a ovde izrazito) nije dobar podatak za opis populacije. Odgovarajućim testiranjem \footnote{Primenjen je Shapiro-Wilkov test za normalnost raspodele. Mogao je biti korišćen i Kolmogorov-Smirnovljev
    test i na ovaj broj opseravcija bi davali isti zaključak, ali zbog preciznosti i brzine rada je on izbegnut.} je pokazano da
    većina polja koja se prati ne ispunjava uslove \footnote{U statistici, prosečna vrednost je merilo populacije onda i samo onda kad veličina podleže tzv normalnoj raspodeli. Njena osobina
je da se tada najveći deo vrednosti koje imamo gomila upravo oko te prosečne i da se odstupanja velike većine tih vrednosti od proseka smatraju zanemarljivim. Takođe, tačka u kojoj je prosečna vrednost deli
    sam uzorak na dve jednake polovina tj. čini raspodelu simetričnom. Sve što od odstupa od toga daje raspodelu koja je asimetrična.} da bi mogao preko ovog podatka da opisujem kretanje vrednosti.

    \textbf{Medijana} $Me$ je odabrana kao mera za opisivanje populacije zbog toga što je matematički dokazano da bez obzira na to kako su raspoređene
    vrednosti, ona se uvek nalazi između sredine i modusa \footnote{Najčešća vrednost koja se javlja}. U našem slučaju, pošto je kretanje veličina takvo da imamo nagli skok
    koji je praćen postupnim opadanjem (tzv. desna asimetričnost u odnosu na sredinu), važi da je prosek nešto veći od medijane. Ta razlika je u ovom slučaju prihvatljivo velika
    i ne utiče na preciznost, a pokazuje praktično istu stvar kao prosečna vrednost tamo gde ona zaista sme da se koristi. Za primer pogledati sliku \ref{fig:korelaciona}

    \begin{figure}[htbp]
        \centering
        \includegraphics[width=\textwidth]{"../rezultati/grafikoni/tipican_mesecni_racun.png"}
        \caption{Prosečan mesečni račun tipičnog korisnika}
        \label{fig:churn-po-paketu}
    \end{figure}

    \section*{Nalazi}

    Poslovni interes je zadržati što je moguće više korisnika. Od ukupno $1191$ ispitanih korisnika \footnote{Nad ovima je tehnički izvodljivo ispitivanje !}
    curenja se uočavaju:

     \begin{figure}[htbp]
        \centering
        \includegraphics[width=\textwidth]{"../rezultati/grafikoni/churn_po_seg.png"}
        \caption{Odlazak korisnika nakon godinu dana po segmentima}
        \label{fig:churn-po-seg}
    \end{figure}


    \begin{table}[htbp]
    \centering
    \caption{Grupe sa najvišim stopama churn-a}
    \label{tab:najvece-churn-stope}
    \small
    \begin{tabular}{l l r}
    \toprule
    \textbf{Kategorija} & \textbf{Grupa} & \textbf{Stopa odlaska (\%)} \\
    \midrule

    Pol
    & Ž
    & 23.06 \\

    Tip paketa
    & Prepaid
    & 29.56 \\

    Kanal prodaje
    & Partner
    & 24.46 \\

    Starosna grupa
    & 30--44
    & 24.95 \\

    Staž
    & $< 1$ god.
    & 32.57 \\

    \bottomrule
    \end{tabular}
    \end{table}

    \subsection*{Povezanost odlaska sa ostalim poljima}
    Pogledati sliku \ref{fig:korelaciona} kako bi se videla jačina povezanosti odlazaka sa drugim veličinama.

    \begin{figure}[htbp]
        \hspace*{-3cm}
        \includegraphics[width=1.5\textwidth]{"../rezultati/grafikoni/korelaciona_mat.png"}
        \caption{Korelaciona matrica odlaska i drugih veličina}
        \label{fig:korelaciona}
    \end{figure}

    \subsection*{Neka istraživačka pitanja}

    \subsubsection*{Da li korisnici odlaze zbog nekog kanala prodaje ?}
    Na osnovu raspoloživih 1191 upisa, verovatnoća da je ovo zaista realan trend je \textbf{manji od $1\%$} što kazuje da apsolutno ne postoji
nikakav prostor za sumnju ovog tipa (videti sliku \ref{fig:ulaz-izlaz})

    \textbf{Napomena}: Vrednost $1$ označava odlazak !

    \subsubsection*{Postoje li obrasci korišćenja sadržaja paketa prema starosnoj grupi ?}

    \begin{table}[htbp]
    \centering
    \caption{Rezultati testa razlika između starosnih grupa}
    \label{tab:kruskal-starosne-grupe}
    \small
    \begin{tabular}{lrrl}
    \toprule
    \textbf{Promenljiva} & \textbf{H vrednost} & \textbf{Verovatnoća slučajnosti} & \textbf{Zaključak} \\
    \midrule
    Potrošnja interneta & 478.490 & $2.19 \times 10^{-103}$ & Statistički značajna razlika \\
    Slanje poruka & 2.105 & 0.551 & Nije statistički značajno \\
    Minuti poziva & 6.007 & 0.111 & Nije statistički značajno \\
    \bottomrule
    \end{tabular}
    \end{table}


    Definitivno vidimo da je potrošnja interneta zaista različita prema pripadnosti starosnoj grupi, dok su ostale veličine
prilično uravnotežene. Naknadne analize su pokazale da je priča preko telefona i slanje poruka prilično uravnotežene među
grupama.

    Zbog praktičnosti u izveštavanju, uzeću grupaciju od 18-29 jer imam najviše njenih pripadnika, ali račun pokazuje
isti obrazac i na drugim. Na slici \ref{fig:1829-pot} se vidi kretanje potrošnje interneta u odnosu na normalnu raspodelu. Vidi se da bez obzira
na neuklapanje u nju, ona pokazuje prilično stabilan obrazac.

    \subsubsection*{Preferiraju li se neki paketi u nekim kanalima prodaje ?}

    Testiranje je ukazalo da je verovatnoća preferiranja nekog paketa u nekom kanalu prodaje \textbf{manja od} $1\%$.
    Prema tome, nema prostora ni za kakvu sumnju u ovom pogledu, što pokazuje grafik \ref{fig:paket-kanal}

     \begin{figure}[htbp]
        %\hspace*{-3cm}
        \centering
        \includegraphics[width=\textwidth]{"../rezultati/grafikoni/paket_kanal.png"}
        \caption{Preferenca u tipu paketa u odnosu na kanal prodaje}
        \label{fig:paket-kanal-prodaje}
    \end{figure}

    \begin{figure}[htbp]
        %\hspace*{-3cm}
        \centering
        \includegraphics[width=\textwidth]{"../rezultati/grafikoni/kanal_tip.png"}
        \caption{Odnos odlaska i kanala prodaje}
        \label{fig:ulaz-izlaz}
    \end{figure}
    \begin{figure}[htbp]
        %\hspace*{-3cm}
        \centering
        \includegraphics[width=\textwidth]{"../rezultati/grafikoni/potrosnja_18-29.png"}
        \caption{Kretanje potrošnje interneta grupacije 18-29 u odnosu na normalnu raspodelu}
        \label{fig:1829-pot}
    \end{figure}

    \begin{figure}[htbp]
        %\hspace*{-3cm}
        \centering
        \includegraphics[width=\textwidth]{"../rezultati/grafikoni/paket_kanal.png"}
        \caption{Preferencije paketa po kanalu prodaje}
        \label{fig:paket-kanal}
    \end{figure}



    \newpage

        \section*{Prikazi važnijih kretanja}
    U nastavku će biti prikazana kretanja važnijih veličina za potrebe analiziranja ovog segmenta tržišta.
    \begin{figure}[htbp]
        \hspace*{-3cm}
        %\centering
        \includegraphics[width=1.4\textwidth]{"../rezultati/grafikoni/churn_po_staz.png"}
        \caption{Odlazak korisnika po stažu provedenom u TS (desno) i po tipu paketa (levo)}
        \label{fig:paket-staz}
    \end{figure}

    \begin{figure}[htbp]
        \hspace*{-2cm}
        %\centering
        \includegraphics[width=1.2\textwidth]{"../rezultati/grafikoni/churn_po_rekl_ocena.png"}
        \caption{Odlazak korisnika po stažu broju reklamacija (desno) i po oceni zadovoljstva (levo)}
        \label{fig:paket-rekl}
    \end{figure}
    \newpage
     \begin{figure}[htbp]
        \hspace*{-3cm}
        %\centering
        \includegraphics[width=1.4\textwidth]{"../rezultati/grafikoni/boxplot_racun.png"}
        \caption{Prikaz kretanja iznosa računa po statusu otišao/ostao (desno) i prema mesečnom računu (levo)}
        \label{fig:paket-ocena}
    \end{figure}

    \newpage
    \section*{Nalazi, ograničenja i preporuke}
    Najveće ograničenje koje nameće skup podataka jeste njegov obim. Cca 1260 upisa nije ni iz bliza dovoljno da bi
    se na objektivan način opisala populacija korisnika koju pokrivamo. Matrica korealcija odlaska sa ostalim veličinama
    nam pokazuje posledičnu vezu tih veličina (kako kretanje jedne utiče na drugu) bez ulaska u razloge zbog kojih
    je ta veza nastala.

    Najveću stopu odlaska uviđam kod Prepaid korisnika. Podaci pokazuju da korisnici koji nisu zadovoljni obično već posle godinu dana raskinu ugovor, a da već posle istekle prve godine, dosta teže
    se odlučuju da menjaju operatera. Samo na osnovu datog skupa podataka ne mogu da zaključim koji su razlozi za takvo ponašanje. Predlažem da se analizira
    strukturu korisnika Prepaid sistema sa ciljem da se utvrde obrasci korišćenja usluge na osnovu čega bi mogle da se donesu neke realne mere za poboljšanje. Smatram da na promenu operatera
    u prvih godinu dana u Biznis i Postpaid sistemu naplate ne možete imati značajan uticaj. Na osnovu datih podataka se samo može oboriti sumnja da je iznos pretplate faktor
    odlučivanja, budući da se one statistički ne razlikuju značajno.

    Sadržaji koji TS nudi u paketima se prema ovim podacima, ne troše ni iz bliza ravnomerno. Primećuje se osetan porast potrošnje internet saobraćaja u grupi od 18-29
    godina koji dalje progresivno opada kako se ide ka sve starijoj populaciji. Ipak, sa aspekta realnog unapređenja, jako malo podataka imam na raspolaganju.
    Rukovodiocima bih zato savetovao sledeće:
    \begin{itemize}
        \item {\textbf{Dugoročno}, vršiti profilisanje internet saobraćaja na mobilnoj mreži prema relevantnim kriterijumima. Njihovo određivanje bih
        sproveo zajednički sa kolegama iz nadležnog sektora za upravljanje infrastrukturom. Preporučujem da se to za početak vrši
        profilisanje po vrsti saobraćaja (streaming, duge sesije preuzimanja ili postavke sadržaja, navigacija i sl). Na osnovu toga
        \textbf{blago personalizovanu} ponudu. Posebno trebalo bi imati u vidu \textbf{iznos pretplate} budući da su pripadnici ove kategorije
        uglavnom izdržavana ili lica sa nešto nižim primanjima. Pored ovoga, savetujem da se ovaj postupak primeni na pakete koji su u ponudi u poslovnicama ukoliko
        za ovo postoje mogućnosti. }
        \item {\textbf{Kratkoročno} gledati da korekcija cene paketa bude \textbf{blago nepodnošljiva} za preduzeće. Naglašavam blago, budući da korisnika u ovoj grupi
        ima najviše, brz odziv na ovakvu vrstu ponude bi mogao za vrlo kratko vreme da povrati izgubljena sredstva. Određivanje opesga podnošljivosti
        zahteva konsultaciju sa odsekom za finansije. Isto važi i za grupaciju od 30-44 godine koja i najviše napušta. Koliko će egzaktno da se vrše ove korekcije dao bih u razmatranje
        sektoru za finansije.}
    \end{itemize}

    Ostali sadržaji (telefonija i poruke) se koriste uobičajeno i za sada ne uočavam neke značajnije promene u tom pogledu.

    Prema broju reklamacija (slika \ref{fig:paket-rekl}), podjednako bih tretirao korisnike sa 3 i sa 7 i više reklamacija. Savetujem da se unutar nadležnog sektora
    formira posebna grupa ljudi koja će da razmatra sve njih bez obzira na to čiji je propust u pitanju. Jednako je važno udobrovoljiti korisnika
    kojem je podjednaka šansa da nešto prigovara kao i da ne prigovara u toku godine (5,6 reklamacija npr) kao i korisnika koji bi prigovarao
    ređe od jednom u četiri meseca (do 3 reklamacije u poslednjih godinu dana). Pri tome, treba imati na umu da sve i da korisnik NIJE
    u pravu, treba korigovati ponudu na neki način da bi se u buduće izbegle takve vrste neprijatnosti.





\end{document}
